

\newpage

%\footline={\hss\folio\hss}
%\pageno=1

\dspace
\subhead
2. Le point de vue tannakien 
\endsubhead

\subhead
2.1
\endsubhead
Soient $k$ un corps de nombres et $S$ un
ensemble fini de places finies de $k$.
Le foncteur $\omega$, d\'eduit du foncteur
gradu\'e $\omega_*$ de 1.1 par oubli de la graduation,
est un foncteur fibre de la cat\'egorie \`a
produit tensoriel $\MT(\scrO_S)$ d\'efinie en
1.6.
Cette cat\'egorie est donc tannakienne, et 
si $G_\omega$ est le sch\'ema en
groupe des automorphismes du
$\otimes$-foncteur $\omega$, le foncteur
$\omega$ induit une \'equivalence de
$\MT(\scrO_S)$ avec la cat\'egorie des
repr\'esentations (lin\'eaires, de dimension
finie) de $G_\omega$.

L'action de $G_\omega$ sur
$\omega(\dbQ(1))=\dbQ$ d\'efinit un morphisme
$$
G_\omega\to \dbG_m.\tag2.1.1
$$
Soit $U_\omega$ son noyau.
Le groupe multiplicatif $\dbG_m$ agit sur le
$\otimes$-foncteur $\omega$, $\lambda\in\dbG_m$
agissant sur $\omega_n$ par multiplication par
$\lambda^n$.
Cette action est une section $\dbG_m\to
G_\omega$ de (2.1.1).
Elle fait de $G_\omega$ un produit semi-direct
$$
G_\omega=\dbG_m\ltimes U_\omega.
\tag2.1.2
$$

Si cela est n\'ecessaire pour \'eviter une
ambigu{\umi}t\'e, on \'ecrira
$G_\omega\left<\MT(\scrO_S)\right>$ et\break
$U_\omega\left<\MT(\scrO_S)\right>$ au lieu de
$G_\omega$ et $U_\omega$.

L'action de $G_\omega$ est fonctorielle en $M$.
Pour tout $M$ dans $\MT(\scrO_S)$, 
elle respecte donc la 
filtration par le poids, index\'ee
par les entiers pairs, de $\omega(M)$ par les
$$
\omega(W_{-2n}M)=\oplusop_{m\Ge n}\omega_m(M).
$$

Le sous-groupe $U_\omega$ agit trivialement
sur $\omega(Q(1))$,  donc sur les
$\omega(\dbQ(n))$, ainsi que sur $\Gr^W(\omega(M))
=\omega(\Gr^W(M))$.
Il en r\'esulte que le sous-groupe
alg\'ebrique $U_\omega\left<M\right>$ de
$\GL(\omega(M))$ image de $U_\omega$ est
unipotent.
L'action de $\dbG_m\Subset G_\omega$ sur
$\omega(M)$ normalise cette image.
L'alg\`ebre de Lie $\gru_\omega\left<M\right>$
de $U_\omega\left<M\right>$ est donc une
sous-alg\`ebre de Lie gradu\'ee de
$\grgl(\omega(M))$.
Puisque l'action de $U_\omega\left<M\right>$
sur $\omega(M)$ respecte la filtration par le
poids et est l'identit\'e sur le gradu\'e
associ\'e, l'alg\`ebre de Lie
$\gru_\omega\left<M\right>$ est \`a
degr\'es $>0$.

Pour $M$ un motif de Tate mixte, notons
$\left<M\right>$ la sous-cat\'egorie
tannakienne de $\MT(\scrO_S)$
{\it engendr\'ee} par $M$, d'objets les
sous-quotients de sommes de $M^{\otimes
p}\otimes (M\upvee)^{\otimes q}$. 
Si $M'$ est dans $\left<M\right>$,
$U_\omega\left<M'\right>$ est un quotient de
$U_\omega\left<M\right>$, et le sch\'ema en
groupe $U_\omega$ est la limite projective
des $U_\omega\left<M\right>$ pour
$\left<M\right>$ de plus en plus grand.
Le sch\'ema en groupe $U_\omega$ est donc
pro-unipotent.

Les alg\`ebres de Lie
$\gru_\omega\left<M\right>$ sont \`a
degr\'e $>0$ et, par 1.6 et 1.7, les groupes
$$
\alignat2
\Ext^1(\dbQ(0),\dbQ(n)) = &\,0 
   &\qquad&\text{si $n\Le 0$},\tag2.1.3\\
&\scrO_S^*\otimes\dbQ &\qquad&\text{si $n=1$},\\
&K_{2n-1}(k)\otimes\dbQ &\qquad&\text{si $n\Ge 2$}.
\endalignat
$$
sont de dimension finie.
On est donc dans la situation consid\'er\'ee
en A.14.
En chaque degr\'e $n$, le syst\`eme projectif
des $(\gru_\omega\left<M\right>)_n$ est
stationnaire et on d\'efinit l'alg\`ebre de
Lie gradu\'ee $\gru_\omega^{\GR}$ comme
\'etant la somme sur $n$ des 
$\lim(\gru_\omega\left<M\right>)_n$.
L'alg\`ebre de $\Lie\,\gru_\omega$ de
$U_\omega$ est la limite projective des
$\gru_\omega\left<M\right>$; c'est le produit
des $(\gru_\omega^{\GR})_n$.
Appliquant A.15, on obtient:

\proclaim{Proposition 2.2}\quad
(i) L'alg\`ebre de Lie gradu\'ee
$\gru_\omega^{\GR}$ est \`a degr\'es $>0$ et est en
chaque degr\'e de dimension finie.
Elle d\'etermine $U_\omega$ par
$$
U_\omega
=\lim\,\exp(\gru_\omega^{\GR}\text{/partie de
degr\'e $\Ge N$}).
$$

\noindent
(ii)
Le foncteur $\omega$ induit une \'equivalence
de $\MT(\scrO_S)$ avec la cat\'egorie des
espaces vectoriels de dimension finie 
gradu\'es, munis d'une
action de $\gru_\omega^{\GR}$ compatible aux
graduations.
\endproclaim

Que les $\Ext^2(\dbQ(0),\dbQ(n))$ soient nuls
admet la traduction suivante.

\proclaim{Proposition 2.3}
L'alg\`ebre de Lie $\gru_\omega^{\GR}$ est libre.
\endproclaim

On en obtient un syst\`eme g\'en\'erateur
homog\`ene libre en relevant une base de 
chaque
$$
\left(\gru_\omega^{\GR}\right)_n^{\ab}=
\text{ dual de (2.1.3)}
\tag2.3.1
$$
dans $(\gru_\omega^{\GR})_n$.

\subhead
2.4 Mise en garde
\endsubhead
On ne dispose pas d'un rel\`evement
canonique de $(\gru_\omega^{\GR})_n^{\ab}$ dans
$(\gru_\omega^{\GR})_n$.
On ne dispose donc pas d'un isomorphisme
canonique de $\gru_\omega^{\GR}$ avec l'alg\`ebre
de Lie librement engendr\'ee par l'espace
vectoriel gradu\'e
$$
(\gru_\omega^{\GR})^{\ab}=\oplus(\Ext^1(\dbQ(0),
\dbQ(n))\upvee\text{ en degr\'e $n$}).
\tag2.4.2
$$

\subhead
2.5
\endsubhead
La cat\'egorie tannakienne $\MT(k)$ est la
r\'eunion filtrante des sous-cat\'egories
pleines $\MT(\scrO_S)$, pour $S$ de plus en
plus grand.
On a donc pour $\MT(k)$
$$
\align
G_\omega\left<\MT(k)\right>:=
\Aut^\otimes(\omega)
&=\lim\,G_\omega\left<\MT(\scrO_S)\right>=\\
&=\dbG_m\ltimes\lim\,U_\omega
\left<\MT(\scrO_S)\right>.
\endalign
$$

\subhead
2.6
\endsubhead
Rappelons quelques d\'efinitions fondant la
g\'eom\'etrie alg\'ebrique dans une
cat\'egorie tannakienne $\scrT$.

La cat\'egorie $\Ind\,\scrT$ des $\Ind$-objets
de $\scrT$ est munie d'un produit tensoriel
associatif commutatif \`a unit\'e h\'erit\'e
de $\scrT$.
Ceci permet de d\'efinir la notion
d'alg\`ebre commutative \`a unit\'e de $\Ind\,\scrT$:
c'est un objet $A$ muni d'un produit $A\otimes
A\to A$ et d'une unit\'e $1\to A$ v\'erifiant
les axiomes usuels.
On d\'efinit la cat\'egorie des
$\scrT$-{\it sch\'emas affines}, aussi
appel\'es sch\'emas affines en $\scrT$, comme
\'etant la duale de celle des alg\`ebres
commutatives \`a unit\'e de $\Ind\,\scrT$.
On note $\Spec(A)$ le $\scrT$-sch\'ema affine
correspondant \`a l'alg\`ebre $A$, et on
appelle $A$ son {\it alg\`ebre affine}.

Un $\scrT$-sch\'ema en groupe affine est un
objet en groupe de la cat\'egorie des
$\scrT$-sch\'emas affines, i.e. le spectre
$G=\Spec(A)$ d'une alg\`ebre de Hopf
commutative $A$ de $\Ind\,\scrT$.
Une action de $G$ sur un objet $M$ de $\scrT$
est une structure de comodule $M\to A\otimes M$.

Pour $M$ un object de $\scrT$, on notera encore
$M$ le $\scrT$-{\it sch\'ema vectoriel}
$\Spec(\Sym^*M\upvee)$, o\`u
$\Sym^*N:=\colim_A\oplusop\limits_{n\Le A}\Sym^n(N)$. 
Un pro-objet $\lim\,M_i$ de $\scrT$ d\'efinit
de m\^{e}me le $\scrT$-{\it sch\'ema provectoriel}
$\Spec(\colim_{i,A}\oplusop_{n\Le A}
\Sym^n(M_i\upvee))$.

Pour un expos\'e plus d\'etaill\'e, nous
renvoyons \`a Deligne (1989) \S5.
Par loc. cit. 5.8, de nombreux \'enonc\'es
connus en g\'eom\'etrie alg\'ebrique restent
valables en $\scrT$-g\'eom\'etrie alg\'ebrique.
Nous les utiliserons librement.

\subhead
2.7
\endsubhead
Rappelons Deligne (1990) qu'on dispose,
dans toute cat\'egorie tannakienne $\scrT$,
d'un $\scrT$-sch\'ema en groupe affine
canonique, le {\it groupe fondamental}
$\pi(\scrT)$ de $\scrT$.
Il agit sur tout objet $M$ de $\scrT$, cette
action est fonctorielle en $M$, et pour tout
foncteur fibre $F$ elle d\'efinit un
isomorphisme 
$$
F(\pi(\scrT))\simover \Aut^\otimes(F).\tag2.7.1
$$

\subhead
2.8
\endsubhead
Certaines des constructions 2.1--2.5 sont
motiviques, i.e. l'image par $\omega$ de
constructions dans $\MT(\scrO_S)$.
D'autres d\'ependent de ce que le foncteur
fibre $\omega$ est gradu\'e.

Notons $G$ le groupe fondamental $\pi(\MT(\scrO_S))$.
Les constructions suivantes sont motiviques:

\medskip\noindent
{\sans (a)}\enspace
le morphisme (2.1.1) de $G_\omega$ dans
$\dbG_m$.
C'est l'image par $\omega$ du morphisme
$$
G\to\dbG_m
\tag2.8.1
$$
donnant l'action du groupe fondamental $G$ de
$\MT(\scrO_S)$ sur $\dbQ(1)$.
Pour donner un sens \`a (2.8.1), observer que
la cat\'egorie des espace vectoriels de
dimension finie est une sous-cat\'egorie de
$\MT(\scrO_S)$, par $V\mapsto V\otimes\dbQ(0)$,
et qu'un sch\'ema en groupe affine peut donc
\^{e}tre consid\'er\'e comme un sch\'ema en
groupe de $\MT(\scrO_S)$.

\medskip\noindent
{\sans (b)}\enspace
le groupe $U_\omega$ est motivique: c'est
l'image par $\omega$ du noyau $U$ de (2.8.1).
De m\^{e}me pour sa pro-alg\`ebre de Lie
$\gru$

\medskip\noindent
{\sans (c)}\enspace
la formule (2.4.2) se rel\`eve en
un isomorphisme de pro-objets de
$\MT(\scrO_S)$
$$
\Lie(U^{\ab})=\prod\limits_{n}
\Ext^1(\dbQ(0),\dbQ(n))\upvee\otimes\dbQ(n).
\tag2.8.2
$$
Dans (2.8.2), la projecton sur le $n^{\ieme}$
facteur donne l'action de $\Lie\,U$ sur les
extensions de $\dbQ(0)$ par $\dbQ(n)$.

Par contre, la structure de produit
semi-direct (2.1.2) et la graduation de
$\gru_\omega$ (i.e., la d\'ecomposition de ce
provectoriel en produit des
$\left(\gru_\omega^{\GR}\right)_n$)
ne sont pas motiviques.  
La graduation de $\gru_\omega^{\ab}$ l'est,
image de (2.8.2) par $\omega$, car
l'action int\'erieure de $G$ sur $\gru^{\ab}$
se factorise par le quotient $\dbG_m$ de $G$.

Si cela est n\'ecessaire pour \'eviter une
ambiguit\'e, on \'ecrira
$G\left<\MT(\scrO_S)\right>$ et
$U\left<\MT(\scrO_S)\right>$ au lieu de $G$
et $U$.

\subhead
2.9
\endsubhead
La r\'ealisation de de Rham $M_{\DR}$ de $M$
dans $\MT(k)$ est un $k$-espace vectoriel, muni
d'une filtration par le poids $W$ image de
celle de $M$, et d'une filtration
d\'ecroissante $F$, la filtration de Hodge.
Ces filtrations sont fonctorielles en $M$,
exactes et compatibles au produit tensoriel.

On a
$$
\dbQ(1)_{\DR}=k\text{ en poids $-2$ et
filtration de Hodge $-1$},
$$
et $\dbQ(n)_{\DR}$ est donc $k$ en poids $-2n$
et filtration de Hodge $-n$.
Pour $M$ quelconque, les filtrations $W$ et
$F$ sont donc oppos\'ees, et d\'efinissent une
graduation de $M_{\DR}$ telle que
$$
\align
W_{-2n}(M_{\DR}) &=\oplusop_{m\Ge n}
  (M_{\DR})_m\\
F^{-n}(M_{\DR}) &=\oplusop_{m\Le n}
  (M_{\DR})_m.
\endalign
$$
On a
$$
\align
(M_{\DR})_n &=\Gr_{-2n}^W(M_{\DR})=
  (\Gr_{-2n}^W(M))_{\DR}\\
&=(\omega_n(M)\otimes\dbQ(n))_{\DR}=
  \omega_n(M)\otimes k.
\endalign
$$
On en d\'eduit la

\proclaim{Proposition 2.10}
Le foncteur fibre $\DR$ sur $\MT(k)$ est
canoniquement isomorphe au foncteur d\'eduit
du foncteur fibre $\omega$ par extension des
scalaires de $\dbQ$ \`a $k$.
\endproclaim

\subhead
2.11
\endsubhead
Pour $C$ une cl\^{o}ture alg\'ebrique de $\dbR$
et $\sigma$ un plongement de $k$ dans $C$, \`a
la th\'eorie de cohomologie
$$
X\text{ sur } k\longmapsto H^*\text{ (espace
topologique $X(C),\dbQ$)}
$$
correspond un foncteur fibre $M\mapsto
M_\sigma$, fonctoriel en $C$.
Parler de ``plongement dans $C$'' plut\^{o}t
que de plongement complexe (i.e. dans $\dbC$)
a l'avantage que la fonctorialit\'e en $C$
fournit un isomorphisme
$M_\sigma\to M_{\bar{\sigma}}$.
Pour un plongement r\'eel, i.e. pour 
$\sigma=\bar{\sigma}$, cet isomorphisme est
une involution de $M_\sigma$, le {\it Frobenius 
r\'eel}.

Un isomorphisme de comparaison canonique d'un
foncteur fibre $\alpha$ vers un foncteur
fibre $\beta$ sera not\'e
$\comp_{\beta,\alpha}$, et son inverse
$\comp_{\alpha,\beta}$.
La relation entre la cohomologie de de Rham
et la cohomologie singuli\`ere fournit
$$
\comp_{\sigma,\DR}\colon\,\,M_{\DR}\otimes_{k,\sigma}
C\simover M_\sigma\otimes_\dbQ C,
\tag2.11.1
$$
fonctoriel en $C$.
Il induit
$$
\comp_{\sigma,\omega}\colon\,\,
M_\omega\otimes C\simover
M_\sigma\otimes C\tag2.11.2
$$
et permet de voir $M_\sigma$ comme une
$\dbQ$-structure sur
$M_{\DR}\otimes_{k,\sigma}C=M_\omega\otimes
C$.
On a
$$
\dbQ(1)_\sigma=2\pi i\dbQ\Subset
C=\dbQ(1)_{\DR}\otimes_{k,\sigma}C
$$
et donc $\dbQ(n)_\sigma=(2\pi i)^n\dbQ\Subset
C$.
Dans ces formules, le choix de la racine
carr\'ee $i$ de $-1$ n'importe pas.

On reconstitue la $\dbQ$-structure $M_\omega$
de l'espace vectoriel gradu\'e $M_{\DR}$ \`a
partir de $M_\sigma$ et $\comp_{\sigma,\DR}$
par
$$
(M_\omega)_n\Subset(M_{\DR})_n\Subset
(M_{\DR}\otimes_{k,\sigma} C)_n\quad 
\est~\frac{1}{(2\pi i)^n}\,\comp_{\DR,\sigma}
\Gr_{-2n}^W(M_\sigma).
\tag2.11.3
$$

Ecrivons $G_\omega=\dbG_m\ltimes U_\omega$
pour $G_\omega\left<\MT(k)\right>$ et sa
d\'ecomposition 2.5, et $\tau$ pour
l'inclusion de $\dbG_m$ dans $G_\omega$.


\proclaim{Proposition 2.12}
Il existe $a_\sigma$ dans $G_\omega(C)$ tel que
$$
M_\sigma=\comp_{\sigma,\omega}(a_\sigma
M_\omega).\tag2.12.1
$$
Si $i$ est une racine carr\'ee de $-1$ dans
$C$, on peut choisir les $a_\sigma$ de la forme
$$
a_\sigma=a_\sigma^0\tau(2\pi i)\tag2.12.2
$$
avec $a_\sigma^0\in U_\omega(C)$ et
$a_{\bar{\sigma}}^0=(a_\sigma^0)^-$, d'o\`u
pour $\sigma$ r\'eel, 
$$
a_\sigma^0\in U_\omega(\dbR)\qquad
\text{(pour $\sigma=\bar{\sigma})$}
\tag2.12.3
$$
\endproclaim

La preuve est la m\^{e}me que celle de
Deligne (1989) 8.10.
Noter que (2.12.1) d\'etermine $a_\sigma$ \`a
multiplication \`a droite pr\`es par un \'el\'ement
de $G_\omega(\dbQ)$.
Choisissons $i$ dans $C$.
La fonctorialit\'e en $C$ montre que si
$a_\sigma$ v'erifie (2.12.1), alors
$\abar_\sigma$ v\'erifie (2.12.1) pour
$\bar{\sigma}$.
Si on impose \`a $a_\sigma$ d'\^{e}tre de la
forme (2.12.2), avec $a_\sigma^0$ dans
$U_\omega(C)$, alors
$\abar_\sigma=\abar_\sigma^0\tau(-2\pi
i)=\abar_\sigma^0\tau(2\pi i)\tau(-1)$ et
$\abar_\sigma^0$ est bien un choix possible
pour $a_{\bar{\sigma}}$.
Pour la v\'erification qu'on peut m\^{e}me
prendre $a_\sigma^0\in U_\omega(\dbR)$ si
$\sigma=\bar{\sigma}$, on renvoie \`a loc.  cit.
Si on impose \`a
$a_\sigma$ d'\^{e}tre de la forme (2.12.2),
son ind\'etermination est r\'eduite \`a
$U_\omega(\dbQ)$.
Si $\sigma$ est r\'eel et qu'on impose
en outre (2.12.3), elle est r\'eduite \`a
$$
\align
U^+(\dbQ):&=U(\dbQ)\CAP\tau(2\pi i)^{-1}
  U(\dbR)\tau(2\pi i)\\
&=\exp\left(\prod \gru_{2n}\right).
\endalign
$$

\subhead
2.13
\endsubhead
Appelons ``syst\`eme de r\'ealisations de Hodge
de Tate mixte'' (sur $k$) la donn\'ee\break
suivante.

\medskip\noindent
{\sans (a)}\enspace
un $k$-espace vectoriel gradu\'e $V_{\DR}$.
La filtration (croissante et index\'ee par les
entiers pairs) $W_{-2n}:=\oplusop_{m\Ge
n}(V_{\DR})_m$ est la {\it filtration par le
poids} et la filtration d\'ecroissante
$F^{-n}=\oplusop_{m\Le n}(V_{\DR})_m$ la {\it
filtration de Hodge}.

\medskip\noindent
{\sans (b)}\enspace
pour tout plongement $\sigma$ de $k$ dans une
cl\^{o}ture alg\'ebrique $C$ de $\dbR$, un
$\dbQ$-espace vectoriel $V_\sigma$ muni d'une
filtration croissante $W$ index\'ee par les entiers
pairs, fonctoriel en $C$.

\medskip\noindent
{\sans (c)}\enspace
pour $\sigma$ comme in (b), un isomorphisme de
comparaison $\comp_{\sigma,\DR}\colon\,
V_{\DR}\otimes_{k,\sigma}C\simover
V_\sigma\otimes C$, respectant la filtration
par le poids et fonctoriel en $C$.
Son inverse est not\'e $\comp_{\DR,\sigma}$.

On suppose ces donn\'ees telles que le
sous-espace
$$
(V_\omega)_n:=\frac{1}{(2\pi
i)^n}\,\comp_{\DR,\sigma}(\Gr_{-2n}^W(V_\sigma))
\Subset (V_{\DR})_n\otimes_{k,\sigma} C
$$
soit contenu dans $(V_{\DR})_n$ et
ind\'ependant de $\sigma$.

La cat\'egorie $\scrR_k^H$, ou simplement
$\scrR^H$, de ces syst\`emes, munie du produit
tensoriel \'evident, est une cat\'egorie
tannakienne, et 2.9 \`a 2.11 fournissent un
$\otimes$-foncteur $\real^H$ de $\MT(k)$ dans
$\scrR_k^H$.

La remarque suivante permettra de v\'erifier
\`a peu de frais le caract\`ere motivique
d'objets qui nous int\'eressent.

\proclaim{Proposition 2.14}
Le foncteur $\real^H\colon\, \MT(k)\to
\scrR_k^H$ est pleinement fid\`ele, d'image
essentielle stable par sous-objet.
\endproclaim

Le point essentiel est que le foncteur
$\real^H$ induit une injection de
$\Ext^1(\dbQ(0),\dbQ(n))$ dans
$\Ext^1(\real^H(\dbQ(0)),\real^H(\dbQ(n)))$.
Pour $n=1$, cette injectivit\'e est celle de 
$$
(\log\,\sigma)\colon\, k^*\otimes\dbQ\to
\prod\limits_{\sigma\colon\, k\to\dbC}
\dbC/2\pi i\dbQ
$$
(et un $\sigma$ suffit \`a l'injectivit\'e).
Pour $n>1$, elle est une cons\'equence de
l'injectivit\'e de l'application r\'egulateur
sur $K_{2n-1}(k)$.

\demo{Preuve}
Le foncteur fibre $\omega$ identifie la
cat\'egorie $\scrR_k^H$ \`a celle des
repr\'esentations d'un groupe alg\'ebrique
$\dbG_m\ltimes U_\omega^H$, et le foncteur
$\real^H$ au foncteur
$\Rep(G_m\ltimes U_\omega)\to\Rep(\dbG_m\ltimes
U_\omega^H)$ d\'efini par un homomorphisme
$$
\dbG_m\ltimes U_\omega^H\to \dbG_m\ltimes U_\omega.
$$
L'injectivit\'e sur $\Ext^1$ assure que ce
morphisme est surjectif, et 2.14 en
r\'esulte.
\enddemo

\subhead
2.15 Variantes
\endsubhead
(i) Puisque $\MT(\scrO_S)$ est une
sous-cat\'egorie pleine de $\MT(k)$ stable par
sous-objet, 2.14 vaut \'egalement pour le
foncteur
$$
\real^H\colon\, \MT(\scrO_S)\to \scrR_k^H
\tag2.15.1
$$

\noindent
(ii)
Supposons donn\'ee une factorisation du
foncteur 2.15.1 par une cat\'egorie
ab\'elienne $\scrR:$
$$
\real^H=F\circ \real^\scrR\colon\,
\MT(\scrO_S)\to\scrR\to\scrR_k^H,
$$
avec $\real^\scrR$ et $F$ exacts et $F$
fid\`ele.
Le foncteur $\real^\scrR$ est alors pleinement
fid\`ele d'image essentielle stable par
sous-objet.
``Pleinement fid\`ele'' r\'esulte formellement
de ``$F$ fid\`ele'' et ``$F\circ\real^\scrR$
pleinement fid\`ele''.
Que $F$ soit exact et fid\`ele implique que
pour $R$ dans $\scrR$, $F$ induit une
injection
$$
\{\text{sous-objets de $R$}\}\longrightarrow
\{\text{sous-objets de $F(R)$}\}
$$
et ``image essentielle stable par
sous-objet'' r\'esulte formellement de la
m\^{e}me propri\'et\'e pour
$F\circ\real^\scrR$.

On peut par exemple prendre pour $\scrR$ la
cat\'egorie $\scrR^{H+\ell}$ suivante de
syst\`emes de r\'ealisations: un objet est la
donn\'ee de (a) (b) (c) comme en 2.13, plus
les donn\'ees $\ell$-adiques suivantes:

\medskip\noindent
{\sans (d)}\enspace
pour $\ell$ un nombre premier et $\kbar$ une
cl\^{o}ture alg\'ebrique de $k$, un
$\dbQ_\ell$-espace vectoriel $V_\ell$, muni
d'une filtration finie croissante $W$ index\'ee par
les entiers pairs, la filtration par le poids,
d\'ependant fonctoriellement de $\kbar$.
La fonctorialit\'e en $\kbar$ doit d\'efinir une
action continue de $\Gal(\kbar/k)$, et on
demande que la repr\'esentation
$\Gr_{-2n}^W(V_\ell)$ de $\Gal(\kbar/k)$ soit
isomorphe \`a une somme de copies de
$\dbQ_\ell(n)$.

\medskip\noindent
{\sans (e)}\enspace
pour $\sigma$ un plongement de $k$ dans une
cl\^{o}ture alg\'ebrique $C$ de $\dbR$, et
pour $\kbar$ la cl\^{o}ture alg\'ebrique de
$k$ dans $C$, un isomorphisme de comparaison
$$
\comp_{\ell,\sigma}\colon\,
V_\sigma\otimes\dbQ_\ell\simover V_\ell
$$
fonctoriel en $C$.

On suppose l'existence d'un syst\`eme de
r\'eseaux $M_{\sigma,\dbZ}\Subset M_\sigma$
tels que les
$$
\comp_{\ell,\sigma}(M_{\sigma,\dbZ}\otimes \dbZ_l)
\Subset V_\ell
$$
soient stables sous $\Gal(\kbar/k)$ et
ind\'ependants de $\sigma$.

Tant dans $\scrR^H$ que dans $\scrR^{H+\ell}$,
les seuls objets simples sont les
r\'ealisations des $\dbQ(n)$.

\subhead
2.16
\endsubhead
Pour $k'$ une extension finie de $k$, on
dispose d'un $\otimes$-foncteur ``extension du
corps de base \`a $k'$'': $M\mapsto M_{(k')}$
de $\MT(k)$ dans $\MT(k')$.
L'existence et les propri\'et\'es du morphisme
``norme'' en $K$-th\'eorie assurent par
application de (0.4) que 
$$
\Ext_{\MT(k)}^1(\dbQ(n),\dbQ(m))\to
\Ext_{\MT(k')}^1(\dbQ(n),\dbQ(m))
\tag2.16.1
$$
est injectif et que, pour $k'/k$ une extension
galoisienne, l'image est form\'ee des
invariants sous Galois:
$$
\Ext_{\MT(k)}^1(\dbQ(n),\dbQ(m))\simover
\Ext_{\MT(k')}^1(\dbQ(n),\dbQ(m))^{\Gal(k'/k)}.
\tag2.16.2
$$

Par les arguments de la preuve de 2.14,
l'injectivit\'e de (2.16.1) implique que le
foncteur $M\mapsto M_{(k')}$ est pleinement
fid\`ele, d'image essentielle stable par
sous-objet.
En particulier, si un motif de Tate mixte
$M'$ sur $k'$ provient d'un motif de Tate
mixte $M$ sur $k$, ce dernier est
d\'etermin\'e \`a isomorphisme unique pr\`es
par $M'$.
Si la place finie $v'$ de $k'$ domine la
place $v$ de $k$, $M'$ est non ramifi\'e en
$v'$ si et seulement si $M$ l'est en $v$.

\subhead
2.17
\endsubhead
Par 2.16, les cat\'egories tannakiennes
$\MT(k)$, pour $k$ variable, forment un
pr\'echamp sur le site \'etale de
$\Spec(\dbQ)$.
On note $\MAT$ le champ associ\'e et on appelle
$\MAT(k)$ la cat\'egorie des motifs
d'Artin-Tate mixte sur $k$.
Explicitons cette d\'efinition.

Soit $k_1$ une extension galoisienne de $k$.
Notons $\scrC(k_1/k)$ la cat\'egorie des
extensions finies $k'$ de $k$ dans lesquelles
peut se plonger $k_1$.
Soit $\MAT(k_1/k)$ la cat\'egorie des sections
de $\MT$ au-dessus de $\scrC(k_1/k)$: la
cat\'egorie des syst\`emes d'objets
$M_{(k')}\in \MT(k')$ pour $k'\in\scrC(k_1/k)$,
et d'isomorphismes $M_{(k')(k'')}\simover
M_{(k'')}$ pour $k'\to k''$ dans
$\scrC(k_1/k)$, v\'erifiant une
compatibilit\'e pour $k'\to k''\to k'''$.
La cat\'egorie $\MAT(k)$ est la limite
inductive des $\MAT(k_1/k)$, pour
$\scrC(k_1/k)$ de plus en plus petit.

La filtration par le poids des objets des
$\MT(k')$ \'etant fonctorielle et compatible
aux extensions de corps de base, elle fournit
une filtration par le poids sur chaque objet
de $\MAT(k_1/k)$, ou de $\MAT(k)$.

La cat\'egorie $\MAT(k)$ h\'erite de $\MT(k')$
de foncteurs fibre $\omega$, $\DR$, $\sigma$
et $\ell$, et d'isomor-\break
phismes de comparaison.

\medskip\noindent
{\it Foncteur fibre $\omega$}:
le foncteur $M=(M_{(k')})_{k'\in\scrC(k_1/k)}
\mapsto M_{(k_1)}$ est une \'equivalence de
$\MAT(k_1/k)$ avec la cat\'egorie des objets de
$\MT(k_1)$ munis d'une action semi-lin\'eaire
de $\Gal(k_1/k)$
(donn\'ee de descente galoisienne).
Cette donn\'ee de descente induit une action
de $\Gal(k_1/k)$ sur
$\omega(M):=\omega(M_{(k_1)})$ et le groupe de
Galois motivique $\Aut(\omega)$ de
$\MAT(k_1/k)$ est le produit semi-direct de
$\Gal(k_1/k)$ par le groupe de Galois
motivique de $\MT(k_1)$.

Soit $\kbar$ une cl\^{o}ture alg\'ebrique de
$k$ et passons \`a la limite sur $k_1\Subset
\kbar$.
On d\'efinit
$$
\omega(M):=\omega(M_{(k_1)})\qquad
\text{pour $k_1\Subset \kbar$ assez grand}.
$$
Ce $\dbQ$-espace vectoriel gradu\'e est muni
d'une action continue de $\Gal(\kbar/k)$.

\medskip\noindent
{\it Foncteur fibre $\DR$}:
d\'efini par descente galoisienne \`a partir
des foncteurs fibres $\DR$ des $\MT(k')$.
Pour $M$ dans $\MAT(k_1/k)$ et
$k'\in\scrC(k_1/k)$, on a fonctoriellement en
$k'$
$$
M_{\DR}\otimes_k k'\simover (M_{(k')})_{\DR},
$$
et l'isomorphisme 2.10 fournit un isomorphisme
de comparaison 
$$
M_{\DR}=(\omega(M_{(k_1)})\otimes
k_1)^{\Gal(k_1/k)}.
$$

\medskip\noindent
{\it Foncteur fibre $\sigma$}:
pour $\sigma$ un plongement de $k$ dans une
cl\^{o}ture alg\'ebrique $C$ de $\dbR$.
C'est
$$
M_\sigma:=(M_{(k')})_{\sigma'}
$$
pour $\sigma$ le compos\'e $k\Subset k'
\,{\overset{\smash\sigma'}\to \longhook\,} C$ et $k'$
assez grand pour que $M_{(k')}$ soit dans
$\MT_{(k')}$.

\medskip\noindent
{\it Foncteur fibre $\ell$}:
il d\'epend du choix d'une cl\^{o}ture
alg\'ebrique $\kbar$ de $k$, et est d\'efini
comme \'etant $(M_{(k')})_\ell$ pour
$k'\Subset\kbar$ assez grand. 

\medskip\noindent
{\it Cas particulier.}
Les objets purement de poids $0$ de $\MT(k')$,
i.e. tels que $\Gr_{-2n}W=0$ pour $n\not=0$,
sont les sommes de copies de $\dbQ(0)$, et
$\omega_0$ identifie leur cat\'egorie \`a
celle des $\dbQ$-espaces vectoriels.
En cons\'equence, pour $\kbar$ une cl\^{o}ture
alg\'ebrique de $k$, le foncteur $\omega$
identifie la cat\'egorie des objets purement
de poids $0$ de $\MAT(k)$ avec celle des
$\dbQ$-espaces vectoriels munis d'une action
continue de $\Gal(\kbar/k)$ (motifs d'Artin).

Soit $k_1$ une extension galoisienne finie de
$k$.
Il r\'esulte de 2.16 que le foncteur $\MT(k)\to
\MAT(k_1/k)$ est pleinement fid\`ele d'image
essentielle stable par sous-objet.
L'image essentielle est aussi stable par
extensions:

\proclaim{Proposition 2.18}
Pour qu'un objet $M$ de $\MAT(k_1/k)$ soit dans
$\MT(k)$, il faut et il suffit que l'action de
$\Gal(k_1/k)$ sur $\omega(M)$ soit triviale.
\endproclaim

Si $M$ est purement de poids $0$, la
proposition r\'esulte du ``cas particulier''
ci-dessus.
Le cas o\`u $M$ est pur, i.e. purement d'un
seul poids, en r\'esulte par tensorisation avec
un $\dbQ(n)$.
Il reste \`a v\'erifier une stabilit\'e par
extensions.

\demo{$1^{\text{ere}}$ preuve}
Une suite spectrale d'Hochschild-Serre
fournit, pour $M$ et $N$ dans $\MT(k)$, des
isomorphismes
$$
\Ext_{\MAT(k_1/k)}^p(M,N)\simover \Ext^p
(M_{(k_1)},N_{(k_1)})^{\Gal(k_1/k)}.
$$

Pour v\'erifier que
$$
\Ext_{\MT(k)}^p(M,N)\to\Ext_{\MAT(k'/k)}^p
(M,N)=\Ext_{\MT(k_1)}^p(M_{(k_1)},
N_{(k_1)})^{\Gal(k_1/k)}
$$
est bijectif pour $p=0,1$ et injectif pour
$p=2$, il suffit par d\'evissage de le
v\'erifier pour $M$ et $N$ de la forme
$\dbQ(m)$ et $\dbQ(n)$.
Le cas $p=2$ r\'esulte de la nullit\'e des
$\Ext^2$, le cas $p=0$ r\'esulte de ce que
$\Hom=0$, sauf pour $m=n$, o\`u
$\End(\dbQ(n))=\dbQ$, et $p=1$ est 2.16.
\enddemo

Plutot que de d\'efinir la suite spectrale
d'Hochschild-Serre requise, on peut la
paraphraser:

\demo{$2^{\text{\`eme}}$ preuve}
Si $M$ est une extension de $\dbQ(n)$ par
$\dbQ(n+r)$, que $M_{(k_1)}$ admette une
donn\'ee de descente de $k_1$ \`a $k$,
respectant la structure d'extension de
$\dbQ(n)$ par $\dbQ(n+r)$, implique que la
classe de $M_{(k_1)}$ dans
$\Ext_{\MT(k_1)}^1(\dbQ(n),\dbQ(n+r))$ est fixe
sous $\Gal(k_1/k)$.
Elle provient donc de la classe d'une
extension $\Mtil$ de $\dbQ(n)$ par
$\dbQ(n+r)$ dans $\MT(k)$.
On v\'erifie, utilisant que
$\Hom_{\MT(k_1)}(\dbQ(n),\dbQ(n+r))=0$, que $M$
est isomorphe \`a l'image de $\Mtil$.

Prouvons par r\'ecurrence sur $r\Ge 0$ que 2.18
est vrai sous l'hypoth\`ese additionnelle que
pour un $n$ on ait $M=W_{-2n}M$ et
$W_{-2(n+r+1)}M=0$.

Le cas o\`u $r=0$ ($M$ pur) a d\'ej\`a \'et\'e
trait\'e.
Le cas $r=1$ se r\'eduit \`a celui d'une
extension de $\dbQ(n)$ par $\dbQ(n+1)$.

Pour $r\Ge 1$, l'hypoth\`ese de r\'ecurrence
assure que $M'=M/W_{-2(n+r)}M$ et
$M''=W_{-2(n+1)}M$ proviennent de $\Mtil'$ et
$\Mtil''$ dans $\MT(k)$.
Ce sont des extensions de la forme:
$\Gr_{-2n}^WM$ par $N$ et $N$ par
$\Gr_{-2(n+r)}^W M$.
Par nullit\'e des $\Ext^2$ dans $\MT(k)$, il
existe dans $\MT(k)$ un objet $\Mtil$,
extension it\'er\'ee de $\Gr_{-2n}^WM$ par $N$
par $\Gr_{-2(n+r)}^WM$, redonnant $\Mtil'$
et $\Mtil''$.
L'image de $\Mtil$ dans $\MAT(k_1/k)$ diff\`ere de
$M$ par une extension de $\Gr_{-2n}^WM$ par
$\Gr_{-2(n+r)}^W M$.
\enddemo

Cette derni\`ere provient d'une extension dans
$\MT(k)$, par r\'eduction au cas d'une
extension de $\dbQ(n)$ par $\dbQ(n+r)$, 
et ceci permet de corriger $\Mtil$ pour que son
image devienne isomorphe \`a $M$.

